\chapter{Background}
% probably talk about what is SpMV anyway as well, and the family of operations
\section[XLS Overview]{What is XLS?}
XLS, standing for Accelerated Hardware Synthesis, is a development toolchain capable of generating both software and hardware designs from code written in either of its languages: XLScc and DSLX. This flexibility aims to address the growing need of specialized hardware to realize performance gains, enabling users to offload software computations to hardware at their choosing \cite{xlsweb}. XLScc is embedded in C++ and DSLX, 
% bad wording: can i say that?
the language of focus, is their domain specific language.

DSLX provides two functional primitives to describe computations: funcs and procs. Funcs are stateless functions that return a value.
\begin{lstlisting}[caption=An example func, captionpos=b, style=mystyle]
fn mul_add(a: u32, b: u32, c: u32) -> {
    let d = a*b + c;
    (d)
}
\end{lstlisting}
Procs, standing for Communicating Sequential Processes, are stateful modules capable of communicating with one another using directional channels. The model of communication is based upon Kahn Process Networks, where IO operations, such as sends and receives, produce tokens that can be used to order operations with respect to one another. 
\begin{lstlisting}[caption=An example proc. join() produces a new token useful for ordering IO operations \textit{after} all the tokens in its parameter list, captionpos=b, style=mystyle]
proc mul_add{
    a:      chan<u32> in;
    b:      chan<u32> in;
    c:      chan<u32> in;
    d:      chan<u32> out;
    config(
        a: chan<u32> in, b: chan<u32> in,
        c: chan<u32> in, d: chan<u32> out
    ) {(a, b, c, d)}
    init {()}
    next (state: ()){
        let (tok1, av) = recv(token(), a);
        let (tok2, bv) = recv(token(), b);
        let (tok3, cv) = recv(token(), c);
        let tok = join(tok1, tok2, tok3);
        send(tok, d, av*bv + cv);
    }
}
\end{lstlisting}
DSLX generates hardware by lowering the encoded logic into SystemVerilog modules. The pipelines and state machines generated can be additionally configured through compilation flags.

A further introduction to DSLX programming can be found on the XLS website.

% do i talk about XLS inspirations (Handel-C) and/or compare XLS to existing programming languages?
\section[Sparse Matrix Computations]{Sparse Matrix Computations} % should i just combine this with hisparse, i am litearlly just going to say SpMV means sparse matrix which means many zeros lol
% i  think i  do do this section as especially talking about compressed sparse row and compressed sparse column, that will segue into Hisparse architecture overview (where I explain hisparse streaming format from how to make it from CSR, and the streaming benefit it gives compared to CSR)
Sparse matrix computations are a family of matrix operations where many elements of the input matrix are zeros. The operations are often denoted in shorthand with "Sp" highlighting which matrix is sparse in the notation, thus SpMSpM and SpMV indicates a sparse matrix-matrix and sparse matrix-vector operation respectively. As the zeros of the sparse matrix can be trivially ignored when computing the result of the operation, traditional dense algorithms waste both compute resources and useful memory bandwidth when applied to sparse workloads. By instead tailoring the software and/or hardware design for these sparse workloads, performance exceeding a dense computation can be achieved.

% CSR, CSC, COO
Often the storage format of the matrix is optimized to address its sparsity. Compressed Sparse Row (CSR) is a common sparse matrix storage format where a traditional two dimensional array representing the matrix is replaced with three arrays: a value array containing the non-zero matrix elements, a column array containing the column indicies of the elements, and a monotonically increasing row array encoding the row indices of the elements. Unlike the column array, the row array is not the direct row indices of the elements.
% bad wording: below (maybe?)
In other words, it is not the same length as the value array. Instead, the values of the row array represent the index ranges of the value array that reside on a particular row. 
% perhaps figure, perhaps not. Though, i could use a figure to make a worked example with my spmv1.json, showing the changes that HiSparse's format provides and establishing familiarity with the reader.
\begin{figure}[h]
\centering\includegraphics[alt={sample image},width=10cm]{spmv1csr.PNG} 
\caption[Example CSR Matrix]{Example matrix encoded in CSR format}\label{spmv1csr}
\end{figure}
A row array of the form \lstinline|[0, 2, 2, ...]| implies that within the value array, elements 0 and 1 reside on the first row, no elements reside on the second row, and so on. Each successive value in the row array represents the end of the index range of the next row in the matrix. As the name implies, this row array is ideally compressed compared to a simpler direct encoding of the indicies when the input matrix holds many empty rows. Its mirror form, Compressed Sparse Column (CSC), performs the same compression on the column indices of the matrix. Finally, simply storing the value array, and the column and row array in their non-compressed direct encoding results in the Coordinate (COO) storage format.

% is ther innovation on both the storage format as well as the multiplication?

% Existing sparse architectures (nvidia 4:2, ExTensor, GPU vs FPGA?)? Perhaps end with HiSparse, "a SpMV arch used as a case study in this askdoasda" lol
Many current workloads involve sparse matrix computations, and correspondingly there are many options for sparse accelerators within academic and commercial spaces; Introduced in their Ampere architecture, Nvidia's Tensor Cores support full throughput on 2:4 sparse matrices (2:4 implies that for every 4 successive matrix elements, 2 are zeros)\cite{nvidia24}. ExTensor, a generalized sparse tensor accelerator, uses intersections to determine not only when to perform operations between scalars, but also whether to operate on larger granularities such as entire tensors \cite{extensor}. HiSparse, the SpMV accelerator used as a case study for this work, implements its logic in Vitis HLS to show modern high level synthesis capabilities \cite{hisparse}. A further analysis of the HiSparse architecture is provided in the following section.

\section[HiSparse Architecure Overview]{HiSparse Architecture Overview}
HiSparse is a modular FPGA based SpMV accelerator written entirely in Vitis HLS. Expressing dynamic operations such as payload arbitration alongside typical latency insensitive pipeline synthesis, it broadens the idea of what is possible with high level synthesis and likewise serves as an appropriate case study for identifying the capabilities of a hardware focused language. There are slight variations between the integral and floating point version of the architecture, and this overview will focus on the integral version. However, much of the architecture remains the same for the floating point version.
\subsection{High Level Design}
HiSparse is a multi-kernel architecture capable of spanning across chip die boundaries to optimize the frequency of the sythensized design. The number of kernels is flexible and is typically sized with respect to the target matrix workload and FPGA resources. During the SpMV computation, kernels are synchronized with one another on a per row basis, while internal kernel modules can be synchronized in a finer granularity multiple times per row. Synchronization occurs through the communication of tokens that replace data values within the pipeline, thus the propogation of the tokens across the pipeline limits the latency of the design, and the synchronization overhead is amortized through larger matrix computations. Figure~\ref{splitkerneldiagram} shows an example three kernel architecture.

% perhaps figure of the high level split kernel architecture
\begin{figure}[H]
\centering\includegraphics[alt={sample image},width=10cm]{splitkernel.png} 
% option argument of caption suppress citation hyperlink in table of contents by using \CiteNolink rather than \cite
\caption[Split Kernel Diagram~\CiteNolink{hisparse}.]{Example three kernel HiSparse architecture~\cite{hisparse}.}\label{splitkerneldiagram}
\end{figure}

Within the kernel lies one or more clusters, serving as the smallest practically complete computation units; Clusters are capable of requesting matrix data, caching input vector data, computing output vector updates and streaming the final row results out. The internal cluster datapath is comprised of multiple streams running in parallel, each stream computing the results of one or more output vector elements. Streams are numbered successively and the stream number corresponds to the row a stream will iterate over modulo the number of streams. Thus if there are N streams, the first stream will iterate over the first row of the matrix, then it will iterate over the N+1 row of the matrix, and so on. All clusters have the same number of streams, and the streams of a cluster are grouped into structures called channels. Figure~\ref{clusterdiagram} highlights the modules each cluster contains to perform its computations.
% figure of the cluster architecture
\begin{figure}[H]
\centering\includegraphics[alt={sample image},width=10cm]{cluster.PNG} 
\caption[Cluster Diagram~\CiteNolink{hisparse}.]{Cluster diagram~\cite{hisparse}.}\label{clusterdiagram}
\end{figure}

Clusters have localized storage to retain portions of both the input vector and the final output vector. This storage avoids potentially expensive memory requests while processing matrix streams. The storage is split into banks, with each of the cluster's streams owning a particular bank of the cluster's input and output vector storage. All clusters require random access to the same input vector, thus the input vector is duplicated across all clusters and their stream banks. When streaming the final output vector, clusters iterate over their streams' output banks to reconstruct the result.

The size of these banks and the number of streams within a cluster limits the size of the matrix a cluster can process without needing to refresh its storage. The maximum number of columns a cluster can process is equal to the size of a stream's input vector bank multiplied by the number of streams within a cluster. In constrast, the maximum number of rows a cluster can process is equal to the size of a stream's output vector bank multiplied by both the number of streams within a cluster and the number of clusters in the entire design. This discrepancy is due to the input vector being duplicated across all clusters and their stream banks. 

% all clusters have the same number of streams (aka the pack size)
% the input matrix num rows must be divisible by the total number of streams -> ensures matrix loader wont try to load a packed payload type one where only a subset of the PACK_SIZE streams has valid data
% the input matrix num cols must be divisible by the pack size -> ensures the vector loader wont try to grab a packed vector payload when in reality the vector only has valid data for a subset of the PACK_SIZE elements

%hbm channels are really just an analog to pack_size and dont contribute anything to the high level discussion, only to the modified storage format
\subsection{Custom Sparse Matrix Format and Memory Layout}
HiSparse introduces a modified version of the CSR sparse matrix format to address its difficulties in streaming data. To reconstruct the row indicies of a CSR encoded matrix, the row array must be traversed separately from the data and column arrays, taking note of the index ranges of specific rows. Should the matrix hold many empty rows in succession, the reading of the data and column values must be paused until the traversal of the row array has "caught up" to the next valid row. This potentially irregular access pattern hurts streaming throughput.

% figure of this pausing of the row array?
HiSparse replaces the functionality of the row array with tokens in the data and column index arrays to communicate row changes. A token consists of a skip row indicator (represented as a -1 in the column index array) paired with an integer in the data array representing the number of rows to skip. The skip row indicator ensures downstream modules do not treat the data and column pair as valid matrix information, and the integer enables a single token to skip an arbitrary number of rows, rather than the element by element traversal necessary in the CSR row array.

\begin{figure}[H]
\centering\includegraphics[alt={sample image},width=10cm]{spmv1cpsrexample.PNG} 
\caption[Hisparse Matrix Format.]{The previous CSR example expressed in HiSparse's custom format, using two streams and vector and output banks of size 2.}\label{cpsrexample}
\end{figure}
A data and column pair comprises a single payload for a single stream. As mentioned in the previous section, clusters are comprised of one or more streams grouped into channels. Thus, a single channel payload is comprised of one or more stream payloads. These packed channel payloads represent the data type that is read from memory. Consequently, they impose restrictions on the shape of the input matrix; The number of rows of the input matrix must be cleanly divisible by the total number of streams in the entire design, and the number of columns of the input matrix must be cleanly divisible by the number of streams in a cluster. Both requirements ensure there are no memory out of bounds accesses when all of the clusters simultaneously request the packed channel payloads of matrix and vector data at any moment. Vector data is packed in a similar manner to matrix stream data, however, there is no use of tokens as the vector is one dimensional.

% introduce this as how the matrix is prepared in memory PROBLEM -> THIS IS NOT THE ORDER THAT HISPARSE USES HOWEVER THIS WOULD PRODUCE THE SAME RESULT....
The input matrix undergoes multiple preprocessing steps to prepare it for iteration. It must be padded to ensure the divisibility requirements are met. The matrix is also split into equally sized partitions, where the dimensions of the partitions are equal to the maximal number of rows and columns a cluster can process. These partition dimensions must also meet the same divisibility requirements imposed on the input matrix. Finally, after converting the input matrix into the custom format, per-partition metadata is stored in the memory addresses immediately preceding the input matrix. This metadata informs the matrix loading modules of where a given partition is located in memory and the length of each stream within the partition.

\subsection{Key Modules} % better name ?